\section{Discussion\label{sec:discussion}}
The research questions examine where architectural signal resides and how a mapping model can use it. RQ1 compares file representations before message passing, RQ2 examines file, member, and folder context and compares the folder gains with a directory-based rule, and RQ3 examines how predictor choice and aggregation direction affect the contribution of dependency neighbors. This section discusses these findings together.

\parg{Dependency information in file representations.}
Dependency information can support mapping before message passing, through the graph position encoded by Node2Vec or through dependency-guided adaptation of a code encoder. The advantage of these representations over Word2Vec, and of fine-tuned UniXcoder over its zero-shot counterpart, suggests that dependency information is useful beyond what the evaluated lexical and pretrained code representations provide. The representations also differ in training, vector size, and source content, so these comparisons are between representation pipelines and not between isolated information sources.

The usefulness of this signal varies with the relation between dependency structure and the reference architecture. Node2Vec’s micro-F1 is lower on the three least homophilic systems, SH-3D, HDF, and T.Mates (Table~\ref{tab:datasets}), with a median of $58$ compared with $72$ on the remaining seven. Node2Vec learns dependency proximity from files that co-occur on random walks, without using the GT mapping. When dependencies frequently cross module boundaries, files from different modules may co-occur on walks and receive similar vectors, making the modules harder to distinguish.

The advantage of fine-tuned UniXcoder over Node2Vec on SH-3D and T.Mates suggests that code content may provide complementary evidence in these mixed neighborhoods. HDF does not follow this pattern, so low homophily alone does not explain their relative performance; HDF is also the smallest system, and every representation performs comparatively poorly on it. Homophily is used here to interpret the observed pattern, not to select a representation, since computing it requires the reference mapping.
 
Encoding dependencies in the initial representation and aggregating dependency neighbors during mapping are two uses of the same information source, and the second does not necessarily replace the first. Word2Vec and zero-shot UniXcoder gain more from neighbors than Node2Vec and fine-tuned UniXcoder do, yet fine-tuned UniXcoder still has the highest median macro-F1 for every predictor. Conversely, useful dependency information in a representation does not guarantee a benefit from further aggregation: RQ3 shows that this benefit depends on both aggregation direction and predictor choice.

\parg{Folder organization and architectural exceptions.}
The gains from adding folder nodes in RQ2 show that directory organization helps reproduce the reference mapping. However, extending a directory assignment is different from identifying files that depart from it. The directory analysis separates these cases by checking, on the files that a directory-based baseline gets wrong, which incorrect predictions of FM the folder nodes correct and which correct predictions they lose (Section~\ref{sec:directory_errors}).

Extending a directory assignment is different from identifying files that depart from it. Comparing FM and FFM on the directory baseline's errors helps assess this distinction, but does not isolate architectural exceptions from errors caused by missing local supervision.

Lucene, HDF, and A.UML have pure directories under the reference mapping, so their baseline errors are consistent with directories lacking labeled examples, although the analysis does not verify this directly. Compared with FM, folder nodes improve accuracy on these files, but FFM recovers only a minority on A.UML. Pure directories therefore do not ensure that folder nodes recover assignments where local supervision is missing.

Bash and HDC have mixed directories. On the baseline's errors, FFM is similar to FM on Bash and less accurate on HDC. JabRef illustrates the distinction between overall gains and performance on baseline errors: adding folder context raises overall accuracy while lowering accuracy on those errors.

Therefore, a tool intended to identify architectural exceptions should also be evaluated on files whose module differs from the majority module of their directory. Recovering assignments in directories without labeled examples is useful, but does not establish that a model can identify conflicts between directory structure and architecture.

RQ2 uses HGT to model the node and relation types of the heterogeneous variants. Although HGT gains the least from neighbor access in RQ3, its median macro-F1 gap to GAT is only 1.2 to 2.4 points, far smaller than the folder gains in RQ2. This suggests that predictor choice alone is unlikely to explain the folder gains, but whether they persist with GAT on the heterogeneous variants remains untested.

\parg{Dependency neighborhood processing and aggregation direction.}
The benefit of dependency neighbors depends on how the predictor aggregates them. HGT gains least, while GAT achieves the highest median macro-F1 with every representation under undirected aggregation, although differences between predictors are small relative to variation across systems. GAT and HGT differ in their attention formulations and parameterizations, so the comparison does not isolate why GAT performs better. These findings concern the evaluated file-level predictors, not the folder- and member-aware configurations or the usefulness of dependency types in general.

The preferred aggregation direction appears to depend more on the system than on the predictor or representation, although the pattern is not uniform across metrics. This suggests that neither direction is consistently preferable across systems. Agreement across configurations and models describes this consistency rather than providing independent evidence, since the comparisons share systems and splits.

Undirected aggregation adds reverse edges, so the comparison does not isolate propagation direction from neighborhood size. A reversed-edges-only setting would help distinguish these effects and remains for future work. With Node2Vec, directed aggregation lowers median macro-F1 for all three predictors. Since its vectors were learned on an undirected graph, this may reflect a mismatch between the neighborhoods encoded during initialization and those used during prediction, but this explanation remains untested.

\parg{Connectivity and neighbor information.}
Connectivity describes how much neighboring context is available, measured here by average node degree, while homophily describes whether linked files share a module. Few neighbors limit the context available to shallow message passing, whereas Node2Vec can encode more distant context through random walks. Limited connectivity may therefore constrain the benefit of additional aggregation even when the initial structural representation remains useful.

Several results are consistent with this interpretation. With Node2Vec and undirected aggregation, no predictor changes significantly on Bash and HDC, the two lowest-degree systems. In FM and FFM, dependencies between member nodes require two layers to reach file predictions; we therefore examined whether the macro-F1 change from one to two layers relates to average degree. Exploratory Kendall’s $\tau$ analyses~\cite{kendall} show positive associations in FFM across all three representations, but in FM mainly with Word2Vec.

Connectivity may also help explain the effect of aggregation direction. T.Mates and SH-3D combine low homophily with high average degree, gain most from file-level message passing with Node2Vec in RQ2, and favor directed aggregation in RQ3. Bash and HDC favor undirected aggregation, as does HDF, which has low homophily but much lower average degree than T.Mates and SH-3D. Reverse edges may supply missing context in sparse neighborhoods but introduce additional cross-module context in dense, low-homophily neighborhoods. This mechanism remains untested: the ten-system comparison cannot separate connectivity from homophily, system size, or label scarcity.

\parg{Implications for evaluating mapping models.}
Evaluation should establish what the input representation already provides, then isolate access to neighborhood information within a fixed predictor. Comparing a GNN model with an MLP changes both the predictor and its access to dependency structure; enabling and disabling neighborhood information within the same predictor tests whether that access helps.

Micro-F1 and macro-F1 should be read together, since correctly mapping more files does not necessarily improve performance across modules. For tools intended to identify architectural exceptions, aggregate scores should also be accompanied by performance on files whose module differs from the majority module of their directory. These comparisons help distinguish the contributions of representations and neighbor processing from gains that mainly reproduce directory organization.

The intended use also matters. If mapping supports code reorganization, the existing directory structure may be something to reconsider rather than reproduce. Strong performance based on that structure does not establish which information would remain useful when the organization changes. The present study does not test this setting, but motivates evaluating information sources against the intended task rather than treating any representation or predictor as generally preferable.

 